\chapter{Cold-Email Principles}
\label{ch:cold-email-principles}

Chapter~\ref{ch:direct-outreach} gave you one outreach template, aimed specifically at startup founders. This
chapter and the next build a much broader cold-email toolkit, assembled from several independent sources that
converge on remarkably similar principles despite covering different audiences (recruiters, professors, clients,
investors). This chapter covers the \emph{universal} principles; Chapter~\ref{ch:cold-email-templates} covers
ready-to-use templates organized by who you're writing to.

\section{The Five ``Lego Bricks''}

A cold email breaks down into five components, each addressing one common failure mode.

\begin{longtable}{p{0.2\textwidth} p{0.74\textwidth}}
\toprule
\textbf{Brick} & \textbf{What it fixes} \\
\midrule
\endhead
1. Identify the right person & Don't email \texttt{info@company.com} — find whoever actually owns the problem
you solve (a design lead for design roles, a data lead for data roles, and so on). Search LinkedIn, the
company's Team/Leadership pages, or professional directories. \\
2. Begin with why you're reaching out & Don't open with your own backstory. Lead with something relevant to
\emph{their} world: a product launch, a LinkedIn post, a hiring announcement — establish relevance before your
qualifications. \\
3. Make a small request & Don't ask for a job, referral, or meeting in the first message — that's proposing
marriage on a first date. Ask something that costs little to answer: ``Could I ask a quick two-minute
question?'' or ``Could you point me to whoever handles this?'' \\
4. Use meaningful personalization & A generic compliment (``I admire your work'') could be sent to hundreds of
people. Refer to a specific moment, observation, or problem instead — something that demonstrates you actually
looked. \\
5. Follow up professionally & Wait 4--6 days, then send one polite, low-pressure reminder. The tone should
remind, not pressure. \\
\bottomrule
\end{longtable}

\subsection{The Biggest Mistake: Making It About Yourself}

\begin{longtable}{p{0.48\textwidth} p{0.46\textwidth}}
\toprule
\textbf{Self-focused} & \textbf{Recipient-focused} \\
\midrule
\endhead
``I am looking for an opportunity at your company.'' & ``I noticed your team is expanding its data
capabilities.'' \\
``I have always been passionate about AI.'' & ``I noticed your team is building AI-powered products.'' \\
``Please review my resume and consider me for a position.'' & ``I've worked on a related problem and would be
happy to share a relevant idea.'' \\
``I would be grateful if you could give me a job.'' & ``Would you be the right person to ask about this
area?'' \\
\bottomrule
\end{longtable}

\begin{keytakeaway}
The formula: right person + specific reason + relevant connection + small request + polite follow-up. The
strongest emails make the recipient understand three things fast: why you contacted them, why it's relevant to
their work, and how they can respond with minimal effort.
\end{keytakeaway}

\section{Warm Introductions Beat Cold Email}

Before writing a single cold email, check whether you can get a \textbf{warm introduction} instead — someone who
already knows the recipient connecting you or recommending the conversation. One source cites warm introductions
converting at 2--3$\times$ the rate of ordinary cold email (the speaker's own stated experience, not a guaranteed
multiplier). Look through: friends and friends of friends; current and former colleagues; LinkedIn connections;
university alumni networks; alumni of companies you've worked at. Cold email remains useful when a warm path
genuinely doesn't exist — for sales, recruiting, partnerships, user research, or founder-to-founder advice.

\section{Build the Funnel Before You Send Anything}

Rather than picking an arbitrary volume (``100 emails a week''), start from the outcome you want and work
backward. One illustrative B2B funnel:

\begin{longtable}{p{0.4\textwidth} p{0.18\textwidth} p{0.32\textwidth}}
\toprule
\textbf{Funnel stage} & \textbf{Assumed rate} & \textbf{Required quantity} \\
\midrule
\endhead
Emails sent $\rightarrow$ opens & 50\% & 800 emails $\rightarrow$ 400 opens \\
Opens $\rightarrow$ responses & 10\% & 400 opens $\rightarrow$ 40 responses \\
Responses $\rightarrow$ completed demos & 25\% & 40 responses $\rightarrow$ 10 demos \\
Demos $\rightarrow$ customers & 10\% & 10 demos $\rightarrow$ 1 customer \\
\bottomrule
\end{longtable}

Under these assumptions, roughly 800 emails are needed per customer. The point isn't the specific numbers — it's
that every stage loses prospects, so your starting volume has to account for that, \emph{and} the funnel tells
you exactly where to focus: few opens points at targeting/subject lines; opens-but-no-replies points at message
relevance; replies-but-no-meetings points at your CTA; meetings-but-no-sales points at the offer itself.

\begin{caution}
This is the \emph{third} distinct funnel model in this guide, alongside the Outreach Equation's illustrative
rates in Chapter~\ref{ch:outreach-engineering} and the 800$\to$40$\to$4$\to$1 funnel in
Chapter~\ref{ch:direct-outreach}. All three are hypothetical, come from different sources and contexts, and are
not meant to be averaged or reconciled — each illustrates the same underlying point (a funnel is multiplication)
with different assumed rates. Track your \emph{own} numbers rather than expecting any of the three to describe
your results.
\end{caution}

\textbf{Start manually.} Before automating anything, you don't yet know which subject lines work, which
prospects respond, or which messages land. Manual outreach teaches you that directly; automate only once you
understand the process and can consistently produce results.

\section{Targeting Beats Clever Copy}

Better targeting is the single highest-leverage way to improve open rates — a perfectly written email cannot
turn someone with no need for what you offer into an interested prospect. 100 carefully researched emails to
companies with a clear reason to care will usually beat 1,000 generic ones. Practical details that compound with
good targeting: use your personal name as the sender name, not just a company name; keep subject lines short,
relevant, and natural-sounding (``Hey, quick question,'' ``Can I get your advice?'') rather than looking
artificially optimized; email is generally reported as outperforming LinkedIn messages for cold outreach, though
this varies by recipient — some people are flooded with LinkedIn messages, others rarely get them.

\section{Personalize Beyond the Name: Uncommon Commonalities}

Inserting \texttt{[First Name]} is not personalization. The strongest personalization finds an \textbf{uncommon
commonality} — a specific connection between you and the recipient that isn't shared by everyone. Sharing a
gender or nationality rarely counts; discovering you attended the same university and took classes in the same
building does. Other examples: having used the same product for years; a shared professional community; having
faced a similar challenge; a specific feature or product change you genuinely appreciate.

\begin{example}
\bad \textbf{Generic:} ``Hey Aaron, love what you're doing at Creative Market.''

\good \textbf{Specific:} ``I'm a huge Creative Market fan ever since you launched the photos category. It's now
my go-to resource.''
\end{example}

Find these by actually reviewing the person's background, the company's product, and its history — then use the
connection naturally in the opening. Replacing a company name in a template is not the same thing.

\section{Keep It Short, Be Human, Establish Credibility, Focus on the Reader}

Four more principles that recur across every source reviewed for this guide:

\begin{description}
  \item[Keep it short] Recipients often read on a phone, between meetings. A long email gets archived or
  forgotten; a short one gets answered now. Short does not mean vague — keep the context that makes the
  proposal credible, cut everything else.
  \item[Write like a human] Open with ``Hey [Name]'' where appropriate rather than ``Dear Mr.\ X, I hope this
  correspondence finds you well.'' Write the message in your actual editor with one specific person in mind,
  not a template with placeholders. Read it aloud afterward — if something would feel awkward to say, rewrite
  it.
  \item[Establish credibility] From whichever source genuinely applies: a respected affiliation, relevant
  education, past employment, named customers or results, a mutual connection, or demonstrated understanding of
  the recipient's product or problem. The test is always ``why should I believe this person can deliver,'' not
  a list of credentials for their own sake.
  \item[Make it about the reader] Explain the problem \emph{they} have and how you can help, rather than
  listing your own background and hoping it lands. One concrete technique: ask your existing users (or, for a
  job search, people who've seen your work) how they'd describe what you do — their phrasing often communicates
  your value more clearly than your own internal language does.
\end{description}

\section{One Clear Call to Action}

Every email needs exactly one specific, low-effort next step, placed in its own sentence near the end so a
recipient scanning to the bottom can find it immediately. ``Please review our 30-page document and give detailed
feedback'' asks too much; ``Would you be open to a 10-minute conversation about how your team currently handles
this?'' asks something answerable on the spot. If an email tries to get a meeting \emph{and} an introduction
\emph{and} feedback all at once, the recipient has too many decisions and typically acts on none of them.

\section{Learning from Real Examples}

\subsection{Three Emails That Failed}

\begin{description}
  \item[Wrong product for the company] A well-written pitch for a physical-shipping speed tool, sent to a
  company that sold downloadable digital assets. \textbf{Lesson:} verify the prospect has a plausible need
  before investing in personalized copy — no amount of writing quality fixes a targeting mismatch.
  \item[Vague ask, no research] ``Your website and content look interesting — do you collaborate with companies
  to make extra money?'' with an opt-out line at the end. \textbf{Lesson:} be specific about the opportunity; a
  built-in opt-out instruction makes a message feel like a mass campaign.
  \item[All about the sender] Paragraph after paragraph of ``I am... I have... I can...'' with no explanation of
  why the recipient should care. \textbf{Lesson:} your background only persuades when it's tied directly to a
  problem the recipient plausibly has.
\end{description}

\subsection{Three Emails That Worked}

\begin{description}
  \item[The uncommon-commonality email] Subject line referencing a shared university sports team; opened with a
  genuine shared-campus detail; established credibility (former Google employees); respected the recipient's
  time (explained the document had already had a first pass); asked for one specific thing (an email address to
  send a document to).
  \item[The ``special invitation'' template] Mostly templated (only the name and a unique signup link varied),
  but felt personal through informal language, a sense of exclusivity (``one of our first sellers''), and an
  invitation to shape the product. \textbf{Lesson:} templates aren't inherently impersonal — the problem is
  templates that feel generic or irrelevant, not the act of templating itself.
  \item[The reframed smoothie pitch] Original version: ``help our startup hit its Demo Day goals'' (about the
  sender). Revised version: ``we'll throw your whole office a free smoothie party'' (about the recipient).
  \textbf{Lesson:} a good offer gives the recipient a reason to engage that doesn't depend on them caring about
  your goals.
\end{description}

\section{Freelance and Service-Business Outreach: Skill Comes First}

A distinct but closely related framework, aimed at freelancers pitching clients, adds several ideas worth
carrying into any outreach context.

\subsection{The Entitled Freelancer, and Three Concentric Circles}

The source warns against becoming an ``Entitled Freelancer'' — someone who expects clients to hand over work
despite lacking the core skill to deliver it well. Knowing a tool (Photoshop, Canva, Premiere Pro) is not the
same as knowing the underlying skill (visual design, composition, video editing). Picture competition as three
concentric circles: \textbf{Circle 1} is the broad field of freelancers emailing, many without strong skills or
communication; \textbf{Circle 2} is the smaller set whose subject line, communication, and professionalism get
the email opened; \textbf{Circle 3} is the smallest set — people who combine good communication \emph{with} the
actual skill the client needs. Communication and professionalism get you noticed; skill is what gets you hired,
retained, and referred.

\begin{keytakeaway}
The architect-vs-artisan distinction: an artisan executes a design that's already specified; an architect
decides what should be built and why. Freelancer A says ``I know Photoshop and Canva and can create graphics.''
Freelancer B says ``I noticed your recent posts have inconsistent visual hierarchy — I prepared an alternative
carousel that makes the main message easier to scan.'' B is demonstrating judgment, not listing tools — the same
distinction Chapter~\ref{ch:foundations} draws for engineers.
\end{keytakeaway}

\subsection{Personalize Around the Client's Problem, Not Just Their Biography}

The two-plumbers analogy: one plumber fixes the immediate leak and quotes a price. The second notices the pipe
itself is old and failure-prone, recommends replacing it, and offers a modest discount — and gets chosen, because
they diagnosed the underlying problem rather than just the symptom. Applied to outreach: ``I am a video editor,
let me know if you have work'' is the first plumber. ``I noticed your recent videos contain detailed
explanations — I prepared an alternative edit of a short segment showing how tighter pacing could make it easier
to follow'' is the second. The client may not have recognized the problem themselves; naming it specifically and
demonstrating a fix is far stronger than a generic skills pitch — but only when the observation is genuinely
specific, not a speculative guess dressed up as insight.

\subsection{Introduce Yourself Briefly, Then Show Rather Than Tell}

If you have experience, state your profession, relevant work, and one concrete example the recipient can
inspect. If you're a beginner, say so plainly rather than inventing experience — lead with personal projects,
practice work, or a small portfolio instead. Either way, a short, specific work sample built \emph{for this
recipient} (a re-edited clip, a redesigned graphic) demonstrates value far better than a claim.

\subsection{A Low-Friction Call to Action}

Make the request small enough to say yes to immediately — ``All I need is 300 seconds of your time'' reframes
five minutes as something that sounds easy to spare. Ask for one specific thing: a short call, feedback on a
sample, or an introduction to the right person on their team. Stay polite rather than demanding — you're asking
for attention, not owed it.

\subsection{Credibility Through Genuine Referrals and Shared Interests}

Mentioning a real mutual contact gives the recipient a way to verify you beyond your own claims — only invoke a
connection that's genuine, and never imply an endorsement that wasn't given. An authentic shared interest (a
hobby, a shared alma mater, a common professional community) can make an interaction feel more human, but it
should supplement the professional proposal, not replace it — don't spend half the email on a hobby and barely
mention the work you can do.

\subsection{There Is No One Perfect Template}

The recommended shape is subject line $\to$ personalized opening $\to$ brief introduction $\to$ evidence of
value $\to$ one call to action $\to$ an optional, genuine personal note — but the source is explicit that forcing
every email into identical mechanical rules can make a whole genre of outreach start to look the same. Some
effective emails are informal and use humor; others are formal. What actually determines results is whether the
work sample is good and the request is reasonable — polish supports substance, it doesn't substitute for it.

\subsection{Who to Email, and Email Etiquette}

As with internship outreach (Chapter~\ref{ch:cold-email-templates}), make HR your last resort when a more
relevant contact exists — a content lead, creative lead, or marketing manager usually understands the actual
work better. A short, practical etiquette checklist worth running before every send:

\begin{multicols}{2}
\begin{itemize}[nosep]
  \item Use a professional-looking email address
  \item Set up a clean signature (name + one or two relevant links)
  \item Use Reply All only when everyone included needs to stay informed
  \item Use hyperlinked text instead of pasting raw long URLs
  \item Attach files \emph{before} filling in the recipient's address
  \item Format for readability — short paragraphs, selective bold, not an unbroken wall of text
  \item Follow up after 2--3 days if there's no response, briefly and without pressure
\end{itemize}
\end{multicols}

\subsection{Ankur Warikoo's Additions}

Three more points worth carrying forward from a related source: write a \textbf{curiosity-driven subject line}
rather than a generic one (``looking for a job'' never earns a second glance); keep the language simple enough
that \textbf{a fifth-grader could follow it} — the weight should come from the message, not the vocabulary;
state your purpose and your actual ask within the first few lines, using an explicit ``headline'' sentence to
grab attention immediately. The underlying point behind all three: the deep research into the recipient's
specific context — the ``emotional labor'' of truly understanding their situation before writing a word — is
precisely the effort most people skip, which is exactly why doing it is what differentiates a reply-worthy email
from the hundred that get ignored.

\section{Follow Up Without Becoming Annoying}

Two to four manual follow-ups is reasonable; give a few days between each rather than sending in rapid
succession, and make each one add something (new information, a different angle) rather than just restating the
original. Stop if the recipient declines or opts out — don't assume you're owed an explanation, and don't let
persistence curdle into pressure. If you genuinely have nothing new to add, it's fine to let a thread go quiet
and revisit the relationship in a few months rather than keep emailing into silence.

